% Part: sets-functions-relations
% Chapter: size-of-sets-complete

\documentclass[../../../include/open-logic-chapter]{subfiles}

\begin{document}

\olchapter{sfr}{siz}{Gedhene Himpunan}

\begin{editorial}
Bab iki ngrembug enumerasi, himpunan kang bisa dietung, lan kang ora bisa dietung.
Sawetara perangan duwe rong versi: versi kang luwih dhasar,
kang nganggep enumerasi minangka dhaptar utawa surjeksi saka $\PosInt$;
lan versi kang luwih abstrak, kang netepake enumerasi minangka bijeksi karo $\Nat$.
\end{editorial}

\olimport{introduction}

\olimport{enumerability}

\olimport{zig-zag}

\olimport{pairing}

\olimport{pairing-alt}

\olimport{non-enumerability}

\olimport{reduction}

\olimport{equinumerous-sets}

\olimport{comparing-size}

\olimport{schroder-bernstein}

\begin{editorial}
Perangan \olref[sfr][siz][enm-alt]{sec},
\olref[sfr][siz][nen-alt]{sec}, lan \olref[sfr][siz][red-alt]{sec} ing ngisor iki
iku versi alternatif saka \olref[sfr][siz][enm]{sec},
\olref[sfr][siz][nen]{sec}, lan \olref[sfr][siz][red]{sec} kang ditulis Tim
Button kanggo buku Open Set Theory. Versi iki rada luwih
maju lan nggunakake definisi enumerabilitas kang beda, kang luwih
cocog ing teori himpunan: bijeksi karo $\Nat$ utawa
perangan wiwitane, tinimbang bisa didhaptar utawa dadi jangkauan
fungsi surjektif saka $\PosInt$.
\end{editorial}

\olimport{enumerability-alt}
\olimport{non-enumerability-alt}
\olimport{reduction-alt}

\OLEndChapterHook

\end{document}
